\section{Context}
\label{sec:context}

The classical approach to the design of recommendation algorithms
strongly relies on prediciting users' preference signals (e.g., likes or ratings)
~\cite{bobadilla2013recommender,DBLP:journals/csur/ZangerleB23,DBLP:journals/inffus/PanPCC26}.
Such a prediction method can leverage user and/or item similarities \cite{DBLP:journals/csur/SilvaMH24}, 
(nonnegative) matrix factorization \cite{DBLP:journals/tai/AhmadianBRFMJ25},
deep learning algorithms \cite{DBLP:conf/recsys/CovingtonAS16},
graph neural networks \cite{DBLP:journals/tors/GaoZLLQPQCJHL23},
or reinforcement learning \cite{DBLP:journals/csur/AfsarCF23}.

However, by fitting a user's like, swipe and scroll actions, 
these algorithms may create radicalization~\cite{gauthier2026political} 
and affective polarization~\cite{tornberg2022digital},
overemphasize instinctive preferences while dismissing thoughtful judgments~\cite{kahneman2011thinking},
and exacerbate rabbit holes that endanger mental health~\cite{ahmed2024social}.
Such considerations are unfortunately too often deprioritized by academic research.
For instance, the survey~\cite{DBLP:journals/ijon/AljunidHHASA25}
reads ``the main objective of RS is to increase the company's sales.''
Given the recent judicial rulings in the US and the EU about ``addictive social media designs''
that include a qualification of some recommendation systems as ``dark patterns'' 
\cite{eu-tiktok-preliminary,reuters-dutch-meta,instagram-youtube-landmark},
the algorithms that result from this line of research could be considered
to run the risk of producing potentially illegal algorithms.

Conversely, increasing attention has been given to 
``trustworthy recommender systems''~\cite{DBLP:journals/tors/GeLFTLXLXZ25}
with greater focus on privacy~\cite{DBLP:journals/ejis/LiU12}, 
explainability~\cite{DBLP:journals/ftir/ZhangC20}, 
fairness~\cite{DBLP:conf/sigir/DoU22,DBLP:journals/tois/WangM00M23}, 
diversity~\cite{wu2024result}, 
robustness~\cite{DBLP:journals/csur/ZhangCSWTSC26}
and user control~\cite{DBLP:journals/jiis/AtasFEPTU21}.
Our solution \Veche{} arguably scores brilliantly in all these dimensions.
It naturally provides high explainability, 
strong fairness between followed users,
robustness to non-followed users,
and user control through follow decisions and hyperparameter selections.
By upgrading it with Tillé's pivot algorithm~\CITE{}, 
we can additionally provide diversity.
Finally, Section~\ref{sec:interoperability} discusses 
how a decentralization of \Veche{} can provide privacy for following users,
and how homomorphic encryption~\cite{rivest1978data} 
and differential privacy~\cite{DBLP:conf/tamc/Dwork08}
could yield privacy for recommendation senders.

To construct \Veche{}, we built upon inspiring prior work on AI governance.
We especially leverage the concept of simulating a vote 
between the stakeholders' algorithmic representatives introduced by
WeBuildAI~\cite{DBLP:journals/pacmhci/LeeKKKYCSNLPP19} for food donation,
as well as by~\cite{DBLP:conf/aaai/NoothigattuGADR18} and~\cite{DBLP:journals/ai/FreedmanBSDC20}
for autonomous car dilemmas and kidney donation.
Inspired by the more careful psychological analysis 
of human preferences of~\cite{DBLP:conf/aaai/Hoang19,DBLP:conf/hci/LechiakhM23},
\Veche{} also provides a clear distinction 
between a user's preferences for what the recommendations they receive,
and their preferences for what they want their followers to be recommended,
Finally, it leverages alternative platforms that welcome alternative recommendation systems,
in particular BlueSky~\cite{kleppmann2024bluesky} and Tournesol~\cite{hoang2025tournesol}.
